이 문제의 제목은 KGUPC 2026을 개최한 알고리즘 동아리 Solutio(솔루티오)의 슬로건이다.

자연수 $n$에 대하여 함수 $f(n)=n^n$을 생각해 보자. 
예를 들어, $f(1)=1,\, f(2)=4,\, f(3)=27,\, f(4)=256$이므로, $n$이 커질수록 $f(n)$의 값은 매우 빠르게 증가한다는 것을 알 수 있다. 
$O$ 기호는 알고리즘의 복잡도를 표현할 때 사용되는 기호로, Solutio는 ``매우 복잡한 인생을 쉽게 만들어보자''라는 의미를 담아 해당 슬로건을 사용하고 있다.

$f(n)$의 값을 직접 계산해 보던 기룡이는 숫자가 너무 커져 계산하기 어려워지자, 대신 $f(n)$을 $1\,000\,000\,007$로 나눈 나머지를 구하기로 했다.

어떤 정수 $x$를 양의 정수 $m$으로 나눈 나머지를 $x\bmod m$으로 나타낸다. 또한 정수 $a, b$와 양의 정수 $m$에 대하여 다음 성질이 항상 성립한다.
$$(a\times b) \bmod m = \big((a\bmod m)\times(b\bmod m)\big)\bmod m$$

예를 들어, $(17\times 23)\bmod 5 = 391\bmod 5 = 1$ 이고, $\big((17\bmod 5)\times(23\bmod 5)\big) \bmod 5 = (2\times 3)\bmod 5 = 6 \bmod 5 = 1$ 이므로 두 결과가 같다.

이 성질을 이용하면 매우 큰 수를 직접 계산하지 않고도, 곱셈을 수행할 때마다 $m$으로 나눈 나머지만을 유지하며 계산할 수 있다.

자연수 $n$이 주어질 때, $n^n$을 $1\,000\,000\,007(=10^9+7)$로 나눈 나머지를 구해 보자.